\section*{第 \ref{ch:g7:volcanoes} 章——火山}

\begin{solution}{exo:g7:volcanoes:1}
(a) \omterm{def:g7:volcanoes:vent}{火山口}；(b) \omterm{def:g7:volcanoes:magma}{岩浆}；(c) \omterm{def:g7:volcanoes:vent}{火山口通道}；(d) \omterm{def:g4:volcanoes-earthquakes-first-look:eruption}{熔岩}；(e) \omterm{def:g7:volcanoes:magma-chamber}{岩浆房}。
\end{solution}

\begin{solution}{exo:g7:volcanoes:2}
(a) \omterm{def:g7:volcanoes:eruption-types}{溢流式喷发}。(b) \omterm{def:g7:volcanoes:eruption-types}{爆炸式喷发}。(c) \omterm{def:g7:volcanoes:eruption-types}{溢流式喷发}：\omterm{def:g4:volcanoes-earthquakes-first-look:eruption}{熔岩}喷泉
喷出发光的熔滴，熔滴落下后流动；\omterm{def:g7:volcanoes:magma}{岩浆}并未被喷成\omterm{def:g4:volcanoes-earthquakes-first-look:eruption}{火山灰}。
(d) \omterm{def:g7:volcanoes:eruption-types}{爆炸式喷发}：这是\omterm{def:g7:volcanoes:pyroclastic-flow}{火山碎屑流}。
\end{solution}

\begin{solution}{exo:g7:volcanoes:3}
(a) 正确。(b) 错误：\omterm{def:g4:volcanoes-earthquakes-first-look:volcano}{火山}气体中含量最多的是水蒸气；
其次是二氧化碳。(c) 错误：\omterm{def:g7:volcanoes:pyroclastic-flow}{火山碎屑流}是一种
由热气体、\omterm{def:g4:volcanoes-earthquakes-first-look:eruption}{火山灰}和\omterm{def:g3:stones-pebbles-sand:rock}{岩石}碎块组成的雪崩。(d) 错误：它们位于
太平洋板块内部，距最近\omterm{def:g7:plate-tectonics:plate-boundary}{板块边界} 3200~千米以上，位于一个
\omterm{def:g7:volcanoes:hotspot}{热点}之上。
\end{solution}

\begin{solution}{exo:g7:volcanoes:4}
(a) 约 1000 至 1200~度。(b) 约 600 至 800~度。(c) 靠近
\omterm{def:g7:volcanoes:vent}{火山口}处，\omterm{def:g4:volcanoes-earthquakes-first-look:eruption}{熔岩}呈黄色（1000 至 1200~度）而非橙色
（800 至 1000~度）：\omterm{def:g4:volcanoes-earthquakes-first-look:eruption}{熔岩}随离开\omterm{def:g7:volcanoes:vent}{火山口}而逐渐冷却。
\end{solution}

\begin{solution}{exo:g7:volcanoes:5}
(a) $838 + 221 = 1059$ 座\omterm{def:g4:volcanoes-earthquakes-first-look:volcano}{火山}。(b) $1059 \div 1214 \approx 0.872$：
约 \textbf{87~\%}。(c) $69 + 18 + 12 + 0.5 = 99.5$，约 100~\%：
这一微小差异来自对每个百分数取整。
\end{solution}

\begin{solution}{exo:g7:volcanoes:6}
(a) $564 \div 660 \approx 0.855$：约 \textbf{85~\%}。(b)
$660 - 564 - 62 - 33 = 1$：一座陡峭的锥形\omterm{def:g4:volcanoes-earthquakes-first-look:volcano}{火山}。(c) 陡峭的锥体由黏稠的\omterm{def:g4:volcanoes-earthquakes-first-look:eruption}{熔岩}和\omterm{def:g4:volcanoes-earthquakes-first-look:eruption}{火山灰}堆成，
这是\omterm{def:g7:volcanoes:eruption-types}{爆炸式喷发}的标志。在环太平洋
\omterm{def:g4:volcanoes-earthquakes-first-look:volcano}{火山}带上，一个板块俯冲于另一个板块之下，我们预料会有很多\omterm{def:g7:volcanoes:eruption-types}{爆炸式喷发}，
例如圣海伦斯\omterm{def:g4:volcanoes-earthquakes-first-look:volcano}{火山}和皮纳图博\omterm{def:g4:volcanoes-earthquakes-first-look:volcano}{火山}的\omterm{def:g4:volcanoes-earthquakes-first-look:eruption}{喷发}。
\end{solution}

\begin{solution}{exo:g7:volcanoes:7}
(a) $1170 - 1140 = 30$~度。(b) $30 \div 10 = 3$~度每
千米。(c) 在\omterm{def:g4:volcanoes-earthquakes-first-look:eruption}{熔岩}隧道中，\omterm{def:g4:volcanoes-earthquakes-first-look:eruption}{熔岩}与空气不接触：
上方已凝固的顶盖像锅盖一样保存热量。在
露天处，其表面立即向空气散失热量。
\end{solution}

\begin{solution}{exo:g7:volcanoes:8}
(a) 考艾岛，至少早 $5.5 - 0.7 = 4.8$ 百万年。(b) 考艾岛：雨水
与海浪\omterm{def:g5:water-shapes-the-land:erosion}{侵蚀}它已有长得多的时间。(c) 大岛，
新\omterm{def:g4:volcanoes-earthquakes-first-look:eruption}{熔岩}仍在喷出的地方。(d) 向西北：考艾岛先被带离\omterm{def:g7:volcanoes:hotspot}{热点}，
现在位于\omterm{def:g7:volcanoes:hotspot}{热点}的西北方。
\end{solution}

\begin{solution}{exo:g7:volcanoes:9}
(a) 几乎所有\omterm{def:g4:volcanoes-earthquakes-first-look:volcano}{火山}都位于\omterm{def:g7:plate-tectonics:plate-boundary}{板块边界}上（列表中为 87~\%）：
离任何边界都很远的\omterm{def:g4:volcanoes-earthquakes-first-look:volcano}{火山}是例外。(b) 在
夏威夷之下的\omterm{def:g6:inside-the-earth:layers}{地幔}深处，\omterm{def:g7:volcanoes:hotspot}{热点}能保持高温数百万年。它
使其上方的\omterm{def:g3:stones-pebbles-sand:rock}{岩石}部分熔化。\omterm{def:g7:volcanoes:magma}{岩浆}比周围的固体\omterm{def:g3:stones-pebbles-sand:rock}{岩石}
轻，穿过板块上升并\omterm{def:g4:volcanoes-earthquakes-first-look:eruption}{喷发}。随着板块移动，新的
\omterm{def:g4:volcanoes-earthquakes-first-look:volcano}{火山}在\omterm{def:g7:volcanoes:hotspot}{热点}上方生成，老的\omterm{def:g4:volcanoes-earthquakes-first-look:volcano}{火山}则被带离。(c)
橙色点（“位于板块内部”）。
\end{solution}

\begin{solution}{exo:g7:volcanoes:10}
第一批\omterm{def:g7:volcanoes:magma}{岩浆}在上升途中已失去大部分气体，好比一瓶
开盖或静置的汽水：\omterm{def:g7:volcanoes:magma}{岩浆}到达地表时，
已无气体将其炸裂，它便缓缓溢出。6 月 12 日和
15 日的\omterm{def:g7:volcanoes:magma}{岩浆}仍保有气体，如同被摇晃的瓶子：近地表
压强骤降，气体一下子形成气泡，气泡
膨胀，把黏稠的\omterm{def:g7:volcanoes:magma}{岩浆}撕裂成\omterm{def:g4:volcanoes-earthquakes-first-look:eruption}{火山灰}。第一批\omterm{def:g7:volcanoes:magma}{岩浆}即是
已失去气体的那一批。
\end{solution}

\begin{solution}{exo:g7:volcanoes:11}
(a) $10 \div 80 = 0.125$~小时，$0.125 \times 60 = 7.5$~分钟。(b)
$10 \div 130 \approx 0.077$~小时，约 $0.077 \times 60 \approx
4.6$~分钟，即约 5~分钟。(c) 这家人需要 20~分钟撤离，而
熔流仅需 5 至 8~分钟到达：看见熔流时已经太晚了。这家人
只有在科学家看到预警迹象时就离开，才能
在\omterm{def:g4:volcanoes-earthquakes-first-look:eruption}{喷发}之前保证安全。
\end{solution}

\begin{solution}{exo:g7:volcanoes:12}
(a) $111 \times 360 = 39\,960$：约 \textbf{40\,000~千米}。(b)
$40\,000 \div 15 \approx 2667$：每天约 2\,670~千米。(c)
$2667 \div 24 \approx 111$：约 \textbf{110~千米每小时}。(d) 由
风带动：细小\omterm{def:g4:volcanoes-earthquakes-first-look:eruption}{火山灰}一旦被\omterm{def:g4:volcanoes-earthquakes-first-look:eruption}{喷发}抛出，
自身便没有速度；它漂浮在高空，随风飘移，同时
缓缓下落。
\end{solution}

\begin{solution}{pb:g7:volcanoes:1}
\textbf{1.} 自 3 月 16 日至 31 日：$31 - 16 = 15$~天；再加上
4 月的 30~天和 5 月的 18~天：$15 + 30 + 18 = 63$~天。

\textbf{2.} $10\,000 \div 62 \approx 161$：平均每天有超过 160 次
小\omterm{def:g4:volcanoes-earthquakes-first-look:earthquake}{地震}。

\textbf{3.} 隆起表明，\omterm{def:g7:volcanoes:magma}{岩浆}已在
山体内部升得很高，并把\omterm{def:g3:stones-pebbles-sand:rock}{岩石}向外推。它还使北侧
变得陡峭且不稳定，随时可能滑动。

\textbf{4.} 蓝色三角形标示一个板块俯冲于另一个板块之下的位置。由此
可预料\omterm{def:g7:volcanoes:magma}{岩浆}黏稠、富含气体，会发生\omterm{def:g7:volcanoes:eruption-types}{爆炸式喷发}，并形成陡峭的锥体（
即层状\omterm{def:g4:volcanoes-earthquakes-first-look:volcano}{火山}）——正是第一张照片所示的情形。

\textbf{5.} $240 \div 60 = 4$~千米每分钟。

\textbf{6.} 隆起的一侧在压力下把\omterm{def:g7:volcanoes:magma}{岩浆}封住，
如同瓶盖。山体一侧滑开后，压强
立即下降；溶于\omterm{def:g7:volcanoes:magma}{岩浆}中的气体一下子变成气泡，
从敞开的一侧冲出——即侧向喷出，那里正是盖子已经脱落之处。

\textbf{7.} $27 \div 480 \approx 0.056$~小时，$0.056 \times 60 \approx
3.4$~分钟。爆炸气浪至少以该速度移动，因此至多
约需 3.4~分钟。

\textbf{8.} $8 \div 130 \approx 0.062$~小时，$0.062 \times 60 \approx
3.7$：约 \textbf{4~分钟}。

\textbf{9.} $24 \div 8.8 \approx 2.7$：烟柱升起的高度约为
珠穆朗玛峰高度的 2.7 倍。

\textbf{10.} \omterm{def:g7:volcanoes:eruption-types}{溢流式喷发}：黏稠的\omterm{def:g4:volcanoes-earthquakes-first-look:eruption}{熔岩}未发生爆炸便涌出，
在\omterm{def:g7:volcanoes:vent}{火山口}上方堆积。与静置的瓶子一样，这批\omterm{def:g7:volcanoes:magma}{岩浆}
已失去大部分气体。

\textbf{11.} \omterm{def:g4:volcanoes-earthquakes-first-look:eruption}{喷发}前，这座山是一个表面光滑、顶部尖锐的锥体，覆盖着
积雪，其下为绿色森林。\omterm{def:g4:volcanoes-earthquakes-first-look:eruption}{喷发}后，山顶与北侧已不复存在：
一个朝向我们敞开的巨大\omterm{def:g7:volcanoes:vent}{火山口}取而代之，而
前方的山坡与\omterm{def:g1:landscapes-around-us:landform}{山谷}已被\omterm{def:g4:volcanoes-earthquakes-first-look:eruption}{火山灰}和碎石染成灰色。

\textbf{12.} $2950 - 2549 = 401$：这座山失去了约
\textbf{401~m} 的高度，约 400~米，正如在四年级时所讲。
\end{solution}
